% GENERATED FILE. Edit canonical lessons, then run npm run generate:books.
% canonical-source-hash: fnv1a64:5792351978cfab45
% canonical-lessons: GE-C500-sammlung, GE-C500-muenze, GE-C500-basteln, GE-C500-raetsel, GE-C500-kreuzwortraetsel

\chapter{Sammeln und rätseln}
\label{ch:ge-sammeln-und-ratseln}
\hlchaptermodality{500}

\begin{chapteropening}
\textbf{By the end of this chapter:} \emph{I can talk about collecting, crafts and puzzles.}

Complete the last lesson of chapter 500: I can talk about collecting, crafts and puzzles.
\end{chapteropening}

\section[die Sammlung]{die Sammlung — collection}

\label{lesson:GE-C500-sammlung}



\begin{quote}
Before the new one: say the German for to crochet, then the German for to embroider.
\end{quote}

\subsection*{You'll want to know: die Sammlung}
\textbf{die Sammlung} — \textquotedblleft{}collection\textquotedblright{}.

\begin{quote}
Mein Opa hat eine große Sammlung von alten Briefmarken.
\end{quote}

\begin{quote}
\emph{My grandpa has a big collection of old stamps.}
\end{quote}

Plural: \emph{die Sammlungen}.

\begin{cousinweb}
\emph{Sammlung} is built on \emph{sammeln} \textquotedblleft{}to collect\textquotedblright{} with \emph{-ung} (EtymDB); \emph{sammeln} is inherited from Middle High German \emph{sameln}.
\end{cousinweb}

\begin{grammarlens}[title={What you've built}]
One more word for telling, explaining and describing at B1.
\end{grammarlens}

\subsection*{Your turn}
\begin{itemize}
\raggedright
  \item \emph{Say it:} \emph{die Sammlung}
  \item \emph{Say it:} \emph{die Sammlung}, once more
  \item \emph{Say it:} a sentence of your own with \emph{die Sammlung}
  \item \emph{Recall:} say the German for to crochet, then the German for to embroider, then say \emph{die Sammlung} again
\end{itemize}

\begin{tcolorbox}[breakable,colback=teal!4,colframe=teal!35!black,title={Before you move on}]
What does die Sammlung mean? (\textquotedblleft{}Collection\textquotedblright{}.) Say it once more.
\end{tcolorbox}
\section[die Münze]{die Münze — coin}

\label{lesson:GE-C500-muenze}



\begin{quote}
Before the new one: say the German for to embroider, then the German for collection.
\end{quote}

\subsection*{You'll want to know: die Münze}
\textbf{die Münze} — \textquotedblleft{}coin\textquotedblright{}.

\begin{quote}
Er sammelt Münzen aus aller Welt.
\end{quote}

\begin{quote}
\emph{He collects coins from all over the world.}
\end{quote}

Plural: \emph{die Münzen}. Also a mint, where coins are made.

\begin{cousinweb}
\emph{Münze} is inherited from Old High German \emph{muniz}, from Proto-Germanic \emph{munitą}, an early borrowing from Latin \emph{monēta} \textquotedblleft{}place for making coins, coined money\textquotedblright{} (EtymDB).
\end{cousinweb}

\begin{grammarlens}[title={What you've built}]
One more word for telling, explaining and describing at B1.
\end{grammarlens}

\subsection*{Your turn}
\begin{itemize}
\raggedright
  \item \emph{Say it:} \emph{die Münze}
  \item \emph{Say it:} \emph{die Münze}, once more
  \item \emph{Say it:} a sentence of your own with \emph{die Münze}
  \item \emph{Recall:} say the German for to embroider, then the German for collection, then say \emph{die Münze} again
\end{itemize}

\begin{tcolorbox}[breakable,colback=teal!4,colframe=teal!35!black,title={Before you move on}]
What does die Münze mean? (\textquotedblleft{}Coin\textquotedblright{}.) Say it once more.
\end{tcolorbox}
\section[basteln]{basteln — to do crafts, to make things by hand}

\label{lesson:GE-C500-basteln}



\begin{quote}
Before the new one: say the German for collection, then the German for coin.
\end{quote}

\subsection*{You'll want to know: basteln}
\textbf{basteln} — \textquotedblleft{}to do crafts, to make things by hand\textquotedblright{}.

\begin{quote}
Mit den Kindern bastle ich Sterne aus Papier.
\end{quote}

\begin{quote}
\emph{I make paper stars with the children.}
\end{quote}

\emph{Ich bastle} — the \emph{e} of \emph{-eln} usually drops.

\begin{cousinweb}
The origin of \emph{basteln} is not settled; EtymDB has no entry for it.
\end{cousinweb}

\begin{grammarlens}[title={What you've built}]
One more word for telling, explaining and describing at B1.
\end{grammarlens}

\subsection*{Your turn}
\begin{itemize}
\raggedright
  \item \emph{Say it:} \emph{basteln}
  \item \emph{Say it:} \emph{basteln}, once more
  \item \emph{Say it:} a sentence of your own with \emph{basteln}
  \item \emph{Recall:} say the German for collection, then the German for coin, then say \emph{basteln} again
\end{itemize}

\begin{tcolorbox}[breakable,colback=teal!4,colframe=teal!35!black,title={Before you move on}]
What does basteln mean? (\textquotedblleft{}To do crafts, to make things by hand\textquotedblright{}.) Say it once more.
\end{tcolorbox}
\section[das Rätsel]{das Rätsel — riddle, puzzle}

\label{lesson:GE-C500-raetsel}



\begin{quote}
Before the new one: say the German for coin, then the German for to do crafts.
\end{quote}

\subsection*{You'll want to know: das Rätsel}
\textbf{das Rätsel} — \textquotedblleft{}riddle, puzzle\textquotedblright{}.

\begin{quote}
Das Kind löst gern Rätsel in der Zeitung.
\end{quote}

\begin{quote}
\emph{The child likes solving puzzles in the newspaper.}
\end{quote}

Plural: \emph{die Rätsel}. Built on \emph{raten} — to guess.

\begin{cousinweb}
\emph{Rätsel} is built on \emph{raten} \textquotedblleft{}to guess, to advise\textquotedblright{}, inherited from Old High German \emph{rātan}, from Proto-Germanic \emph{rēdaną} (EtymDB): something to guess.
\end{cousinweb}

\begin{grammarlens}[title={What you've built}]
One more word for telling, explaining and describing at B1.
\end{grammarlens}

\subsection*{Your turn}
\begin{itemize}
\raggedright
  \item \emph{Say it:} \emph{das Rätsel}
  \item \emph{Say it:} \emph{das Rätsel}, once more
  \item \emph{Say it:} a sentence of your own with \emph{das Rätsel}
  \item \emph{Recall:} say the German for coin, then the German for to do crafts, then say \emph{das Rätsel} again
\end{itemize}

\begin{tcolorbox}[breakable,colback=teal!4,colframe=teal!35!black,title={Before you move on}]
What does das Rätsel mean? (\textquotedblleft{}Riddle, puzzle\textquotedblright{}.) Say it once more.
\end{tcolorbox}
\section[das Kreuzworträtsel]{das Kreuzworträtsel — crossword}

\label{lesson:GE-C500-kreuzwortraetsel}



\begin{quote}
Before the new one: say the German for to do crafts, then the German for riddle.
\end{quote}

\subsection*{You'll want to know: das Kreuzworträtsel}
\textbf{das Kreuzworträtsel} — \textquotedblleft{}crossword\textquotedblright{}.

\begin{quote}
Jeden Sonntag macht mein Vater das Kreuzworträtsel.
\end{quote}

\begin{quote}
\emph{Every Sunday my father does the crossword.}
\end{quote}

Plural: \emph{die Kreuzworträtsel}.

\begin{cousinweb}
\emph{Kreuzworträtsel} is built on \emph{Kreuz} \textquotedblleft{}cross\textquotedblright{}, which came through Old High German \emph{chruzi} from Latin \emph{crux} (EtymDB), \emph{Wort} \textquotedblleft{}word\textquotedblright{} and \emph{Rätsel} \textquotedblleft{}puzzle\textquotedblright{}.
\end{cousinweb}

\begin{grammarlens}[title={What you've built}]
One more word for telling, explaining and describing at B1.
\end{grammarlens}

\subsection*{Your turn}
\begin{itemize}
\raggedright
  \item \emph{Say it:} \emph{das Kreuzworträtsel}
  \item \emph{Say it:} \emph{das Kreuzworträtsel}, once more
  \item \emph{Say it:} a sentence of your own with \emph{das Kreuzworträtsel}
  \item \emph{Recall:} say the German for to do crafts, then the German for riddle, then say \emph{das Kreuzworträtsel} again
\end{itemize}

\begin{tcolorbox}[breakable,colback=teal!4,colframe=teal!35!black,title={Before you move on}]
What does das Kreuzworträtsel mean? (\textquotedblleft{}Crossword\textquotedblright{}.) Say it once more.
\end{tcolorbox}
